\section{结论与启示}
\label{sec:conclusion}

本报告用两个互补的视角考察中国农业的周期性：第一部分的\emph{品种视角}
回答“周期从哪里来、由什么决定、如何相互传导”；
第二部分的\emph{公司视角}回答“在周期中，谁赚到了钱、谁需要什么”。

\subsection{五个结论}

\paragraph{结论一：农业周期的驱动力量已经从需求转向供给与政策。}
在 2024---2025 年的这一轮下行中，中国的粮食与肉类\hl{需求并未出现显著萎缩}：
CPI 食品分项保持温和、水产品消费仍在增长。
价格下跌的主要原因是供给的恢复与产能的过冲：
玉米连续增产、谷物进口大幅收缩，生猪出栏量在产能去化未完成时仍在增长，
白羽肉鸡祖代更新量创历史新高。\emph{这意味着分析农业周期时，
供给端的产能数据（能繁母猪、祖代种鸡、投苗量、棉花种植面积）
比需求端的宏观数据更重要}。

\paragraph{结论二：不同品种的周期长度由“产能调整的生物学时间”决定。}
本报告给出了一个可比较的时间刻度：蛋鸡与肉鸡约 \num{6}---\num{12} 个月、
生猪约 \num{10}---\num{12} 个月、肉羊约 \num{18}---\num{24} 个月、
肉牛 \num{24}---\num{44} 个月以上、多年生作物（苹果、橡胶、糖料）\num{3}---\num{7} 年。
\hl{时间越长的品种，价格调整的幅度越大、拐点的持续性越强}——
肉牛的供给收缩要等到 2027 年之后，而禽类的产能已经完成了一轮出清与反转。

\paragraph{结论三：跨品种的相关性远弱于直觉。}
在第 \ref{sec:corr} 章的统一口径（月度收益率、同一时间窗口）下，
5\% 显著性阈值 $|\rho|>0.24$ 之上只有三个品种：花生 \num{+0.48}、粳米 \num{+0.29}
与豆一 \num{+0.26}（后两者仅略高于阈值）；
玉米 \num{+0.09}、豆粕 \num{+0.13} 等主力饲料原料品种的相关系数均不显著。
\hl{“猪价影响一切农产品”是一种叙事上的便利，而非统计事实}——
真正的共性来自共同的成本（饲料）与共同的宏观环境，
而这些因素在月度频率上已经被品种自身的供需波动淹没。

\paragraph{结论四：上市公司的盈利分布比行业均值更极端。}
按 2026 年半年报口径，\num{104} 家农业上市公司中 \num{46} 家亏损（\num{44.2}\%）、
\num{44} 家半年 ROE 为负，ROE 中位数仅 \num{0.75}\%（未年化）；
\hl{生猪养殖 \num{11} 家公司在 2026 年上半年全部亏损}，
合计亏损 \num{163.8} 亿元（上年同期为合计盈利 \num{154.1} 亿元）。
\hl{周期底部的行业整合不是渐进的，而是在同一时间、以“头部失血、尾部出清”
的方式一次性完成}。
与 2025 年年报（\num{34} 家亏损、ROE 中位数 \num{2.64}\%）相比，
亏损面进一步扩大至 \num{46} 家，归母净利润合计由上年同期的 \num{+270.9} 亿元
转为 \num{-126.1} 亿元，\emph{报表层面的底部比价格层面的底部更晚出现}。

\paragraph{结论五：周期底部的融资需求是刚性的、而且以补血为主。}
按 2026 年半年报口径，\num{14} 家公司同时处于亏损与高负债（$>$ \num{65}\%）状态，
其中 \num{6} 家负债率超过 \num{80}\%；
\num{31} 家公司连续三年未分红。
这 \num{14} 家公司的负债合计 \num{1855} 亿元、净资产合计仅 \num{694} 亿元，
若要把资产负债率压回 \num{65}\% 的警戒线，需要补充的净资产超过 \num{2000} 亿元。
近三年公告中“担保与关联交易”类公告数（\num{3184}）远高于
“对外投资与转型”类（\num{274}），\hl{说明农业上市公司当前的资本运作主题
是维持资产负债表，而不是扩张}。

\subsection{对投资者的启示}

\begin{enumerate}[leftmargin=1.5em, itemsep=2pt]
  \item \textbf{在周期底部，优先选择“供给刚性 + 政策底”的品种。}
        肉牛（进口保障措施、\num{44} 个月产能周期）与蛋鸡（存栏同比下降 \num{6.1}\%）
        是当前供给端最确定的两个方向；而玉米、小麦等政策收储品种
        在成本线附近的波动更适合区间交易而非趋势投资。
  \item \textbf{在同一产业链中，选择需求与价格脱钩的环节。}
        动保（半年 ROE 中位数 \num{2.95}\%、亏损面 \num{11.1}\%、
        资产负债率中位数 \num{30.2}\%）在农业整体下行中保持了正回报，
        海洋捕捞（\num{9.79}\%）则是资源属性带来的另一类“脱钩”。
        \hl{“卖水人”与“资源型”是本轮周期中最有效的两个抗周期逻辑}。
        需要提示的是，宠物食品的半年 ROE 中位数已回落至 \num{4.00}\%，
        高成长赛道同样开始承受竞争压力。
  \item \textbf{对高负债亏损公司保持警觉，而不是抄底}。
        \num{6} 家负债率超过 \num{80}\% 且亏损的公司中，多数需要
        通过股权融资或资产重组才能维持经营，
        \hl{本报告推断这类公司的市值中包含相当比例的“壳价值”与“重整预期”，
        而非盈利能力的现值}——该判断基于市值与盈利的背离，
        仍需公司层面的债务结构与重组进展进一步验证。
  \item \textbf{用公告与产能数据做交叉验证，而不是依赖单一指标。}
        本报告的方法论支持一点：\emph{产能数据（能繁母猪、祖代存栏、投苗量）
        决定方向，公告数据（定增、重组、担保）决定风险，
        财务数据（负债率、ROE）决定谁能等到周期反转}。
\end{enumerate}

\subsection{尚待解决的问题}

本报告在数据上仍有三处限制，也构成后续研究的方向：
\textbf{其一}，由于公开数据源缺少牛、羊、禽、水产的连续现货价格序列，
这些品种的周期判断更多依赖产业统计（存栏、产量、均价）而非高频价格，
无法与生猪期货做严格的月度频率对比；
\textbf{其二}，水产品与部分经济作物的数据粒度只到年度，
限制了周期相位的精确划分；
\textbf{其三}，上市公司画像中的“未来融资需求”是基于财务约束
与公告行为的推断，尚不包含实际的债券到期结构与银行贷款条款等内部信息，
若需精确测算需引入更细的债务数据。

\begin{keybox}[一句话总结]
\normalsize
农业的周期不是一条曲线，而是\hl{一组由生物学时间、政策边界与资产负债表
共同决定的时钟}——\emph{看懂每个品种的时钟现在指在几点，
比预测下个月的价格更有意义}。
\end{keybox}
